% Данный файл распространяетсяя под лицензией CC BY 4.0. Текст лицензии размещён на https://creativecommons.org/licenses/by/4.0/
% (c) Кафедра системного программирования, 2025

\documentclass[a4paper]{article}

\usepackage[a4paper, top=8mm, bottom=8mm, left=8mm, right=8mm]{geometry}

\usepackage{polyglossia}
\setdefaultlanguage[babelshorthands=true]{russian}
\setotherlanguage{english}

\usepackage{fontspec}
\setmainfont{FreeSerif}
\newfontfamily{\russianfonttt}[Scale=0.7]{DejaVuSansMono}

\usepackage[tiny, compact]{titlesec}

\usepackage{titling}
\setlength{\droptitle}{-1cm}
\pretitle{\begin{center}\begin{bfseries}\Large}
\posttitle{\par\end{bfseries}\end{center}}
\preauthor{\begin{center}\normalsize}
\postauthor{\par\end{center}\vspace{-1.8cm}}

\usepackage{hyperref}
\usepackage{bookmark}
\usepackage{csquotes}
% Сюда можно дописывать нужные вам пакеты.

\title{Разработка инфраструктуры бенчмаркинга графовых алгоритмов поиска пути с ограничениями}

\author{Глазунов Иван Сергеевич}

% Дата много места занимает, её лучше не указывать.
\date{}

\begin{document}

\maketitle

\begin{flushright}
    Группа: \emph{23.Б10-мм}

    Кафедра: \emph{Системного программирования}

    % Степень/звание и должность научника мы лучше вас знаем, их можно не указывать.
    Научный руководитель: \emph{Григорьев Семён Вячеславович}

    % А вот про консультанта укажите официальное место работы, с "ООО" или чем-то таким, и должностью желательно по трудовой, не просто "techlead". Выясните у консультанта.
    Консультант: \emph{Белянин Георгий Олегович. ООО "ВК Цифровые Технологии", разработчик}

    Номер семестра практики: \emph{5}
\end{flushright}


\section{Постановка задачи}

% Буквально парой предложений ввести в контекст, обязательно сформулировать цель работы. Если работа делается где-то глубоко внутри проекта, на введение можно потратить больше места, чтобы подвести читателя от сути проекта в целом к тому, что конкретно вам надо было сделать. Только без воды, никакого "в современном мире"!
Работа выполняется в рамках научно-исследовательского проекта по исследованию алгоритмов исполнения запросов с ограничениями на пути,
задаваемыми регулярными (RPQ) и контекстно-свободными (CFPQ) грамматиками.
\textbf{Целью работы является} разработка и реализация инфраструктуры для запуска, тестирования и сравнения алгоритмов поиска путей с ограничениями, а также их оптимизаций.

% Это пояснение актуальности. Тут обязательно надо обрисовать полезность вашей работы для конечного пользователя (и заодно как-то описать, кто такой ваш конечный пользователь --- может, это соседняя команда). Даже если вы делаете работу "глубоко внутри", выясните у консультанта, чем это людям поможет. И опишите, кто выступал в роли "заказчика" вашей работы, т.е. кто и зачем предложил задачу. Можно без конкретных имён, но названия организаций вполне желательны.
Результаты работы направлены на стандартизацию и упрощение процесса сравнения производительности разрабатываемых
алгоритмов\footnote{RPQ-решатель LA'n'EGG RPQ, URL: \url{https://github.com/SparseLinearAlgebra/la-n-egg-rpq} (дата обращения: 04.05.2026)}
\footnote{Single-Source Regular Path Querying, URL: \url{https://arxiv.org/abs/2412.10287} (дата обращения: 04.05.2026)}.
Конечными пользователями такой инфраструктуры являются исследователи и разработчики графовых алгоритмов, которым требуется воспроизводимо запускать эксперименты на общих наборах данных и сопоставимых настройках.

\section{Описание предлагаемого решения}

% Раздел пишется в свободной форме, с минимумом технических подробностей (можно приводить ссылки на более подробные документы).
На текущий момент для сравнения разрабатываемых алгоритмов не хватает единой точки входа,
которая позволяла бы одинаковым образом работать с разными форматами графов, наборами запросов и реализациями алгоритмов.

% Тут надо кратко описать идею решения, спозиционировать его относительно предыдущих наработок\footnote{Если на них надо сослаться, то в подстраничной сноске, URL: \url{http://example.com} (дата обращения: не забываем указывать).} и существующих аналогов (тоже предельно кратко), по возможности обосновать принятые решения.
Предлагаемое решение заключается в разработке библиотеки и консольной утилиты для выполнения запросов с ограничениями на пути и запуска воспроизводимых бенчмарков.
В рамках работы реализована поддержка нескольких форматов входных графов, разбор запросов на подмножестве SPARQL property paths,
единый интерфейс для подключения алгоритмов, сбор статистики запусков и восстановление прогресса после прерывания бенчмарка.

С архитектурой решения можно ознакомиться в репозитории проекта.
Реализация выполнена на языке Rust; для вычислений над разреженными булевыми матрицами используются SuiteSparse:GraphBLAS и алгоритмы библиотеки LAGraph.
Сборка проекта автоматически подготавливает нативные зависимости: GraphBLAS и LAGraph собираются статически, поэтому пользователю не требуется устанавливать их вручную.
Проект также подготовлен к распространению как Rust-библиотека через crates.io и снабжён автоматически публикуемой документацией на docs.rs.

% Вот это важно, опишите кратко, что использовали
% Реализация решения осуществлялась с использованием \dots (тут описать используемые технологии).

\section{Эксперименты}

% Название секции можно поменять на более подходящее вашей работе

% Если есть отзывы команды, консультанта и т.п., хорошо. Если есть только тесты. то хотя бы их тут опишите.

% Если работа подразумевала замеры чего-то (например, производительности),
% приведите тут итоговые результаты и выводы, и дайте ссылку на таблицы с данными.
% Не забывайте про матстат (приводите матожидание и среднеквадратичное отклонение, не просто результат замера).
Целью эксперимента была оценка накладных расходов при запуске LARPQ через библиотеку и CLI Pathrex по сравнению с запуском исходной реализации, а также сравнение с MillenniumDB.

Pathrex показывает конкурентоспособные результаты: полный путь выполнения через библиотеку и CLI даёт около 24\% накладных расходов относительно исходного LARPQ, но при этом остаётся быстрее MillenniumDB примерно в 3.8 раза на простых запросах и в 1.4 раза на сложных запросах.
Таким образом, добавленная инфраструктура бенчмаркинга не скрывает преимущества алгоритма и позволяет воспроизводимо сравнивать реализации на крупных графах.

Подробнее ознакомиться с экспериментами можно в репозитории \url{https://github.com/VanyaGlazunov/pathrex-bench}

\section{Заключение}

В ходе выполнения работы получены следующие результаты:

% ниже списком перечисляете то, что выносится как результат практики.

\begin{itemize}
    \item Спроектирована архитектура библиотеки и утилиты для тестирования алгоритмов поиска путей с ограничениями
    \item Поддержаны форматы входного графа CSV, MatrixMarket и RDF, а также разбор RPQ-запросов на основе SPARQL property paths
    \item Встроены два RPQ-вычислителя: \texttt{LARPQ} и \texttt{RPQ-Matrix}, использующие алгоритмы LAGraph
    \item Реализован запуск бенчмарков с сохранением подробной статистики и восстановлением после прерывания
    \item Добавлены unit-, интеграционные и end-to-end тесты для форматов данных, построения графа, CLI и RPQ-вычислителей
    \item Настроены публикация Docker-образов, распространение пакетов через crates.io и публикация документации на docs.rs
\end{itemize}


Репозиторий: \url{https://github.com/SparseLinearAlgebra/pathrex}.

Документация: \url{https://docs.rs/pathrex/latest/pathrex/}

Решение находится в разработке.
В дальнейшем планируется добавить алгоритмы для контекстно-свободных грамматик и поддержку оптимизаций LA'n'EGG RPQ.

\end{document}
